
\subsubsection{2.1 Generation-Based Hallucination Detection}

The dominant detection paradigm operates at LLM generation time. SelfCheckGPT \citep{Manakul2023SelfCheckGPT} draws multiple stochastic outputs and measures their consistency via NLI, BERTScore, or n-gram overlap, treating inconsistency as a proxy for hallucination probability. Verbalized confidence elicitation \citep{Kadavath2022Language} prompts instruction-tuned models to express uncertainty directly, though reliability varies by model and generation mode and is inaccessible for non-instruction-tuned models. Each of these approaches requires running the LLM at detection time.

\subsubsection{2.2 NLI-Based Factual Consistency}

Maynez et al. [2020] demonstrated that NLI entailment scores correlate with faithfulness in abstractive summarization. SummaC \citep{Laban2022SummaC} systematized this with sentence-level NLI aggregation (SummaCConv), achieving 74.4\% balanced accuracy on summarization consistency benchmarks. TRUE \citep{Honovich2022TRUE} evaluated DeBERTa-NLI as a factual consistency indicator across 11 datasets, establishing cross-task applicability. Neither work reports AUROC on HaluEval hallucination labels under a balanced pair evaluation. FActScore \citep{Min2023FActScore} proposed atomic-fact NLI verification but requires retrieval infrastructure. ORION \citep{Gerner2025ORION} demonstrated F1 = 0.83 on RAGTruth using post-hoc NLI encoders without generation.

The present work differs from these in scope and framing: it evaluates frozen NLI applied directly to pre-existing HaluEval pairs without retrieval or fine-tuning, across all three task types, with AUROC as the primary metric.

\subsubsection{2.3 Positioning}

\begin{table}[htbp]
\centering
\resizebox{\columnwidth}{!}{%
\begin{tabular}{llll}
\toprule
Method & LLM at Detection? & Multi-task & AUROC on HaluEval \\
\midrule
SelfCheckGPT-NLI [Manakul et al., 2023] & Yes (multi-sample) & Yes & 0.48 (dialogue), 0.53 (QA) \\
TRUE [Honovich et al., 2022] & No & Yes & Not reported on HaluEval \\
SummaC [Laban et al., 2022] & No & Summarization only & Not reported on HaluEval \\
ExtrospectiveNLI (this work) & No & Yes (3 tasks) & 0.709 (dialogue), 0.644 (QA) \\
\bottomrule
\end{tabular}
}
\end{table}

Note: SelfCheckGPT was evaluated on base (non-instruction-tuned) Meta-Llama-3-8B in this study, which produces near-uniform stochastic samples. This configuration likely represents a lower bound on SelfCheckGPT performance relative to its intended deployment with instruction-tuned models. The comparison should be interpreted with this in mind.

\subsubsection{2.4 Hallucination Taxonomy}

The commission/omission distinction corresponds to Maynez et al.'s [2020] intrinsic vs. extrinsic categorization of summarization errors. Commission-type hallucinations (fabricated facts, factual substitutions, entity invention) directly contradict the grounding context; omission-type hallucinations (abstractive gaps, paraphrastic reductions, missing information) fail to capture source meaning without explicit contradiction. NLI contradiction detection is structurally aligned with commission-type hallucinations and misaligned with omission-type. SummaC's strong performance on AggreFact --- which contains explicitly inconsistent (commission-type) summaries --- reflects benchmark alignment with the NLI detection mechanism. To our knowledge, this paper provides the first multi-task AUROC-based quantitative characterization of this boundary across dialogue, QA, and summarization.
